\chapter[Validated Adapter Utility Across Robots and Tasks]{Validated Adapter Utility Across Robots and Tasks}\label{chapter:bench-evaluation}

\begin{quote}\itshape
This chapter presents Exp. 3. It tests whether validated adapters remain
useful across robot morphologies, model families, runtime routes, task
settings, and baselines.
\end{quote}

\section[Introduction]{Introduction}
This experiment tests whether the harness and role-separated agent workflow can
produce adapters that satisfy declared contracts and remain useful beyond a
single generation run. It addresses RQ1 and RQ2 by separating adapter synthesis
performance from downstream adapter use.

\section[Background Literature]{Background Literature}
The evaluation is informed by software benchmarks, coding-agent evaluation,
robot task benchmarks, and tool-use evaluation. These areas motivate the need
to measure not only whether code is generated, but whether the resulting tool
layer is reusable and valid under robot-specific constraints.

\subsection[Coding and Program Synthesis Benchmarks]{Coding and Program Synthesis Benchmarks}
Coding benchmarks measure whether models can produce working software
artifacts. AutoAdapter-Bench extends this logic to robot runtime adapters,
where generated code must satisfy morphology, backend, and primitive
constraints.

\subsection[Robot Task and Tool-Use Evaluation]{Robot Task and Tool-Use Evaluation}
Robot task evaluation measures whether an agent can execute behaviours in an
environment. Tool-use evaluation measures whether an agent can call an existing
interface effectively. This chapter keeps these questions separate from the
synthesis of the interface itself.

\subsection[Validation as Evaluation Evidence]{Validation as Evaluation Evidence}
Validation outcomes serve as evidence for contract conformance. They do not by
themselves prove physical deployment success, but they provide a consistent way
to compare models, robots, and specification conditions.

\section[Data Set and Benchmark Design]{Data Set and Benchmark Design}
The benchmark data set consists of robot specifications, runtime profiles,
primitive requirements, model outputs, generated adapter artifacts, validation
reports, and downstream task evidence. These materials allow the experiment to
compare synthesis success and downstream utility across conditions.

\subsection[Robot and Task Coverage]{Robot and Task Coverage}
The benchmark covers variation in robot morphology, control surface, primitive
set, and runtime assumptions. Such variation is necessary because a reusable
adapter claim is only meaningful if it survives embodiment and backend
differences.

\subsection[Model Families and Prompt Conditions]{Model Families and Prompt Conditions}
The benchmark compares model families and generation conditions. This helps
separate failures caused by insufficient model capability from failures caused
by weak specifications, missing context, or inadequate validation feedback.

\subsection[Baselines]{Baselines}
Baselines distinguish direct task scripting from reusable adapter synthesis and
compare harness-supported generation against weaker specification or feedback
conditions.

\section[Aims of Study and Analysis Approach]{Aims of Study and Analysis Approach}
The aim of the experiment is to evaluate whether generated adapters satisfy
declared contracts, whether validated adapters can be reused across tasks, and
how performance varies across robots, models, and baselines.

\subsection[Adapter Synthesis Metrics]{Adapter Synthesis Metrics}
Adapter synthesis is evaluated through validation pass rates, failure
categories, repair outcomes, contract conformance, and evidence of reuse.

\subsection[Downstream Utility Metrics]{Downstream Utility Metrics}
Downstream utility is evaluated by whether task agents can call validated
adapters as tool layers and whether performance differs from direct scripting
or weaker interface conditions.

\section[Benchmark Analysis]{Benchmark Analysis}
The analysis distinguishes three questions: whether an adapter can be
synthesized, whether it passes validation, and whether it remains useful for
downstream task execution.

\subsection[Reusable Adapter Synthesis]{Reusable Adapter Synthesis}
This analysis reports adapter generation success, validation outcomes, repair
patterns, and failure categories. It provides the primary evidence for the
claim that the harness can produce reusable adapters.

\subsection[Downstream Adapter Utility]{Downstream Adapter Utility}
This analysis reports whether validated adapters can serve as tool layers for
downstream agents. It keeps downstream use separate from synthesis so the
thesis does not confuse using an interface with generating one.

\subsection[Use Versus Synthesis]{Use Versus Synthesis}
Some models may use an existing adapter effectively while failing to synthesize
one. This distinction is central to the thesis because tool use and tool
synthesis are different capabilities.

\section[Results]{Results}
The expected results report adapter generation success, validation outcomes,
downstream utility, baseline comparisons, and recurring failure categories
across robot morphologies and model families.

\section[Discussion]{Discussion}
The benchmark evaluation supports claims about reusable adapter synthesis only
within the benchmark and validation boundaries. Its findings motivate the
boundary and runtime integration analysis in Chapter~\ref{chapter:runtime-improvement}.
